% !TeX root = ../main.tex
\chapter*{Dodatek A. Zapytania użyte w ewaluacji}
\addcontentsline{toc}{chapter}{Dodatek A. Zapytania użyte w ewaluacji}
\label{chap:appendix-queries}

% FIRST DRAFT. The exact queries run on 2026-11-23 (read-only, reporting
% views only: no author hash, no text).

Zapytania z tego dodatku dają wyniki z podrozdziałów~\ref{sec:eval-weak-labels} i~\ref{sec:eval-case-study}.
Wszystkie czytają wyłącznie widoki zbioru raportowego, które nie zawierają pseudonimu autora ani treści opinii, a wykonaliśmy je 23 listopada 2026~r. na danych z odtworzenia z 18 października 2026~r.
Zapytanie z listingu~\ref{lst:q-weak} daje macierz z tabeli~\ref{tab:weak-confusion}.

\begin{lstlisting}[float=hbp, language=SQL, caption={Etykiety modelu a rekomendacje autorów recenzji Steam}, label={lst:q-weak}]
SELECT voted_up, sentiment_label, COUNT(*) AS n
FROM `pop-vibe-check.co_reporting_dev.v_events`
WHERE source = 'steam'
GROUP BY 1, 2
ORDER BY 1, 2;
\end{lstlisting}

Zapytanie z listingu~\ref{lst:q-weak-events} daje zestawienie z tabeli~\ref{tab:weak-events}; po zamianie kolumny grupującej na \code{language} daje też porównanie dla języków.

\begin{lstlisting}[float=hbp, language=SQL, caption={Udział rekomendacji i wskaźnik netto modelu w oknach wydarzeń}, label={lst:q-weak-events}]
SELECT
  IFNULL(event_tag, 'general') AS event_tag,
  COUNT(*) AS n,
  ROUND(AVG(IF(voted_up, 1, 0)), 3) AS thumbs_up_share,
  ROUND(SAFE_DIVIDE(COUNTIF(sentiment_label = 'pos')
                    - COUNTIF(sentiment_label = 'neg'), COUNT(*)), 3) AS net_model
FROM `pop-vibe-check.co_reporting_dev.v_events`
WHERE source = 'steam'
GROUP BY 1
ORDER BY 1;
\end{lstlisting}

Zapytanie z listingu~\ref{lst:q-weekly} daje dane do rysunku~\ref{fig:weekly-net}, a zapytanie z listingu~\ref{lst:q-lift} dane do rysunku~\ref{fig:event-lift}.

\begin{lstlisting}[float=hbp, language=SQL, caption={Tygodniowy wskaźnik netto według źródeł}, label={lst:q-weekly}]
SELECT source, period, records, ROUND(net_sentiment, 3) AS net
FROM `pop-vibe-check.co_reporting_dev.v_sentiment_by_period`
WHERE grain = 'week' AND records >= 20
ORDER BY source, period;
\end{lstlisting}

\begin{lstlisting}[float=hbp, language=SQL, caption={Efekt wydarzeń względem tła z tego samego źródła}, label={lst:q-lift}]
SELECT event_tag, source, records, days,
       ROUND(net_sentiment, 3) AS net,
       ROUND(baseline_net_sentiment, 3) AS base_net,
       ROUND(net_lift, 3) AS lift,
       ROUND(volume_multiplier, 2) AS vol_mult
FROM `pop-vibe-check.co_reporting_dev.v_event_lift`
ORDER BY source, event_tag;
\end{lstlisting}
